% Three-model comparison figure: Legacy OS, In-Process Isolation, Confidential VM
% Assumes \usepackage{tikz}, \usepackage{subcaption}, and \usetikzlibrary{positioning, calc, matrix}

\newcommand{\boxwidth}{4.0cm}
\newcommand{\boxheight}{0.8cm}
\pgfmathsetlengthmacro{\appwidth}{\boxwidth/3 - 0.3cm}

\tikzstyle{tcb}=[rectangle,draw,align=center,minimum width=\boxwidth,text width=\boxwidth,minimum height=\boxheight,fill=blue!20]
\tikzstyle{file}=[rectangle,draw,fill=green!20,minimum width=0.8cm,
minimum height=0.5cm,font=\scriptsize]
\tikzstyle{proc}=[rectangle,draw,fill=red!20,minimum width=\appwidth,minimum height=\boxheight,font=\scriptsize]
\tikzstyle{boundary}=[draw=blue,very thick]
\begin{figure}[t]
	\centering
	% -------------------------------------------------------------
	% (a) Legacy OS protection
	% -------------------------------------------------------------
	\begin{subfigure}[t]{0.31\textwidth}
		\centering
		\begin{tikzpicture}[node distance=1.1cm]
			\matrix[column sep=0pt,row sep=5pt] {

				\node[tcb, text height =1.0 cm] (os) {Operating System}; \\
				\node[tcb] (hw) {Hardware (CPU, Memory)};                \\
			};

			\matrix (apps) [above of=os, column sep=7pt,yshift=0.3cm] {
				\node[proc] (p1) {Proc. 1}; &
				\node[proc] (p2) {Proc. 2}; &
				\node[proc] (p3) {Proc. 3};   \\
			};
			% Files below OS
			\node[file, below of=p1,yshift=-.0cm] (f1) {File 1};
			\node[file, below of=p2,yshift=-.0cm] (f2) {File 2};
			\node[file, below of=p3,yshift=-.0cm] (f3) {File 3};

			\node[boundary,fit=(p1)(f1),inner sep=2pt] (b1) {};
			\node[boundary,fit=(p2)(f2),inner sep=2pt] (b2) {};
			\node[boundary,fit=(p3)(f3),inner sep=2pt] (b3) {};
			\node[boundary,fit=(p3)(f3),inner sep=2pt] (b3) {};

			% Kernel boundary
			\draw[boundary] ($(os.north west)+(-0.0,0.15)$) -- ($(os.north east)+(0.0,0.15)$);

			% Legend
			\matrix [draw, above of=apps, yshift=1.5cm, xshift=5cm, column sep=6pt, row sep=2pt] {
				% Trusted
				\node[rectangle,draw,fill=blue!20,minimum width=0.4cm,minimum height=0.4cm] {};  &
				\node[anchor=west] {Trusted};                                                    &
				% Untrusted
				\node[rectangle,draw,fill=red!20,minimum width=0.4cm,minimum height=0.4cm] {};   &
				\node[anchor=west] {Untrusted};                                                  &

				% Object (File)
				\node[rectangle,draw,fill=green!20,minimum width=0.4cm,minimum height=0.4cm] {}; &
				\node[anchor=west] {System Resources};                                             \\

				\node {\begin{tikzpicture}\draw[boundary] (0,0) -- (0.6,0);\end{tikzpicture}};
				                                                                                 &
				\node[anchor=west] {Domain boundary};                                              \\
			};
		\end{tikzpicture}
		\caption{Legacy OS protection introduces protection domains for user processes and kernel.}\label{fig:legacy-prot}

	\end{subfigure}\hfil
	\begin{subfigure}[t]{0.31\textwidth}
		\begin{tikzpicture}[node distance=1.1cm]
			\matrix[column sep=0pt,row sep=5pt] {

				\node[tcb, text height =1.0 cm] (os) {Operating System}; \\
				% \node[tcb] (hyp) {Hypervisor};                           \\
				\node[tcb] (hw) {Hardware (CPU, Memory)};                \\
			};

			\matrix (apps) [above of=os, column sep=7pt,yshift=0.6cm] {
				\node[proc,minimum width=2.5cm, text depth=1.0cm] (p1) {Proc. 1}; &
				\node[proc] (p2) {Proc. 2};                                         \\
			};

			\matrix (comps) [above of=p1,column sep=10pt,yshift=-1.2cm] {
				\node[boundary] (c1) {C$_A$}; &
				\node[boundary] (c2) {C$_B$};   \\
			};
			% Files below OS
			\node[file, below of=c1,yshift=-.2cm] (f1) {File 1};
			\node[file, below of=c2,yshift=-.2cm] (f2) {File 2};
			\node[file, below of=p2,yshift=-.3cm] (f3) {File 3};
			%
			\node[draw=blue,very thick,fit=(p1)(f1)(f2),inner sep=2pt] (b1) {};
			\node[draw=blue,very thick,fit=(p2)(f3),inner sep=2pt] (b2) {};
			% \node[draw,thick,dashed,fit=(p3)(f3),inner sep=2pt] (b3) {};

			% Kernel boundary
			\draw[boundary] ($(os.north west)+(-0.0,0.15)$) -- ($(os.north east)+(0.0,0.15)$);
			% \draw[dashed, thick] ($(p1.north west)+(0.2,0.18)$) -- ($(p1.south west)+(0.2,0.18)$);

		\end{tikzpicture}
		\caption{Hardware-assisted in-process isolation~\cite{erim,hodor,jenny} introduces protection domains for sub-process components.}\label{inprocess-prot}

	\end{subfigure}\hfil
	\begin{subfigure}[t]{0.31\textwidth}
		\begin{tikzpicture}[node distance=1.1cm]
			\matrix[column sep=0pt,row sep=8pt] {

				\node[tcb, text height =1.0 cm] (os) {OS (Guest)}; \\
				\node[tcb,fill=red!20] (hyp) {Hypervisor};         \\
				\node[tcb] (hw) {Hardware (CPU, Memory)};          \\
			};

			\matrix (apps) [above of=os, column sep=7pt,yshift=0.3cm] {
				\node[proc,fill=blue!20] (p1) {Proc. 1}; &
				\node[proc,fill=blue!20] (p2) {Proc. 2}; &
				\node[proc,fill=blue!20] (p3) {Proc. 3};   \\
			};
			% Files below OS
			\node[file, below of=p1,yshift=-.0cm] (f1) {File 1};
			\node[file, below of=p2,yshift=-.0cm] (f2) {File 2};
			\node[file, below of=p3,yshift=-.0cm] (f3) {File 3};

			\node[boundary,fit=(p1)(f1),inner sep=2pt] (b1) {};
			\node[boundary,fit=(p2)(f2),inner sep=2pt] (b2) {};
			\node[boundary,fit=(p3)(f3),inner sep=2pt] (b3) {};

			\node[boundary,fit=(p1)(p2)(p3)(os),inner sep=4pt] (bvm) {};

			% Kernel boundary
			\draw[boundary] ($(os.north west)+(-0.0,0.15)$) -- ($(os.north east)+(0.0,0.15)$);
			% \draw[boundary] ($(hyp.north west)+(-0.0,0.15)$) -- ($(hyp.north east)+(0.0,0.15)$);
			% \draw[dashed, thick] ($(p1.north west)+(0.2,0.18)$) -- ($(p1.south west)+(0.2,0.18)$);

		\end{tikzpicture}
		\caption{Confidential virtual machines~\cite{advanced_micro_device_sev_2022,cca,tdx} establish isolated protection domains for entire VMs.}\label{cvm-prot}
	\end{subfigure}

	\caption{Comparison between legacy OS protection and modern hardware-assisted protection.}\label{fig:protection-models}
\end{figure}
